\chapter[結論]{結論}\label{chapt6}
火山大国である日本において, 
火山ガスの中でも特に毒性の強い二酸化硫黄濃度の監視が重要である. 
本研究は, ラマンDIALの分布観測, 遠隔, 時間連続性に注目し, 
これによる\ce{SO2}観測を提案した. 
本研究の目的はラマンDIALによる\ce{SO2}観測の実現可能性を検討することであり, 
目安として\ce{SO2}濃度設定値に対する±10\%を基準として, 
評価を行った. 


ラマンDIALの統計誤差の
レーザ波長と受信波長構成に対する依存性に着目し, 
最適な波長仕様の検討を行った. 
結果としてレーザ波長\qty{334.6}{nm}を用いて, 
on波長を\ce{O2}アンチストークス散乱波長の\qty{318.0}{nm} 
off波長を\ce{O2}ストークス散乱波長の\qty{353.0}{nm}とする場合が
最も二酸化硫黄測定に最適なラマンDIALの仕様であり, 
信号強度が微弱なアンチストークス散乱も波長選択によって十分に選択可能であることが明らかとなった. 


距離分解能の調整による統計誤差の抑制について検証し, 
\ce{SO2}が\qty{100}{ppm}程度の高濃度領域で要求する測定誤差を達成することがわかった.
干渉ガス濃度の推定による補正の効果についても検証し, 
霧島火山のような\ce{H2S}が極めて高濃度な領域で測定する場合, 
その影響を完全に補正できない可能性があり, 
事前に妥当な\ce{H2S}濃度の推定値を用意するか, 
補助的に\ce{H2S}を測定する手法を併用する必要が示された. 


より現実的な濃度分布を再現するために
プルームモデルをシミュレーションに導入し, 
距離方向に連続的に変化する濃度分布に対して
距離分解能の要求値を検証した. 
距離分解能の限界値である\qty{1.5}{m}でシミュレーションしたところ, 
火口直上の局所的な高濃度分布に対してピークの取りこぼしが発生したが, 
シミュレーション上の点煙源ではなく現実的な面積を持つ火口では十分に濃度分布に追従可能であると評価した. 
また, 運用方法について火口を目視で確認し, 
視線と直行しないようにラマンDIALを設置することでも対応可能である. 


ラマンDIALの受信部に用いる波長選択のための
干渉フィルターについて設計を行い, 
レーザスペクトルや実行吸収断面積について検証を行った. 
ピーク効率\qty{81}{\%}として, 
onチャンネル側は半値幅\qty{9.2}{nm}, 
offチャンネル側は半値幅\qty{10.6}{nm}のフィルター関数でシミュレーションし,  
レーザ波長を十分に遮断しつつ,
目的とする\ce{O2}ラマン散乱信号を選択的に受信可能と評価した. 
